\section{Qual a diferença entre adotar uma solução proprietária como o sistema operacional Windows quando comparado a adoção de uma solução Open source como o sistema operacional Ubuntu? Quais seriam os pontos negativos e positivos de cada abordagem?}

Uma das principais diferenças entre as duas abordagens e uma das primeiras a ser destacadas é o custo. Sistemas proprietários como o Windows exigem aquisição de licença via pagamento, enquanto sistemas Open Source como o Ubuntu podem ser simplesmente baixados gratuitamente. Todavia essa diferença não se resume simplesmente ao preço de aquisição. Uma vez que em geral as pessoas estão mais acostumadas com soluções proprietárias, a adoção de soluções gratuitas também podem gerar custos indiretos, como treinamento de equipe e eventual contratação de suporte especializado, de modo que o custo total de adoção nem sempre acompanha o custo da licença, especialmente para empresas grandes que não são da área de tecnologia (áreas com mais profissionais não familiarizados com soluções Open-Source que e requerem mais treino e adaptação).

Ademais, soluções proprietárias como o Windows não oferecem transparência quanto aos seus mecanismos internos, devido a sua natureza fechada, o que impede a auditoria pública do código e mantém pontos de falha desconhecidos sob segurança por obscuridade, prática já considerada ultrapassada. Já códigos Open Source como o Ubuntu, tem o código-fonte amplamente auditável pela comunidade, permitindo maior conhecimento sobre os processos em execução e de forma geral reduzindo a presença de bloatware. Essa abertura de código também garante flexibilidade, uma vez que o usuário ou a empresa que adota Open Source pode modificar o sistema livremente para adequá-lo às próprias necessidades, algo que soluções proprietárias como o Windows não permitem.

Notar-se-á que ao contrário do que se é comumente pensado, a natureza de Open Source não significa a completa ausência de Bloatware ou telemetria. Destaca-se que há, por exemplo, a coleta de telemetria no Ubuntu, nominalmente os pacotes "Popularity Contest"~\cite{popularityContest} (desabilitado por padrão) e o "Ubuntu Report"~\cite{ubuntuReport} (que na época era opt-out, vindo selecionado na instalação por padrão, hoje foi substituido pelo "Ubuntu Insights"~\cite{ubuntuInsights}, que por sua vez é opt-in). Todavia, a coleta de telemetria em soluções Open Source é geralmente mais transparente, com a possibilidade de auditoria do código-fonte e geralmente opt-in ou ao menos com opção de desativação, enquanto soluções proprietárias como o Windows não permitem tal auditoria e podem estar ocultas ao usuário.

Outro ponto relevante é a dependência de fornecedor. A adoção do Windows tende a atrelar a empresa a todo um ecossistema Microsoft (Active Directory, Office, .NET, entre outros), criando um ciclo vicioso no qual as únicas soluções disponíveis são as do fornecedor e/ou parceiros. Soluções Open Source, por seu direito de fork, reduzem essa dependência, embora as múltiplas distribuições possam gerar fragmentação e exigir mais cuidado na escolha da distro adequada e confusão sobre qual fork exotérica obscura usar.

Outra diferença está em quem responde pelo funcionamento do sistema. Ao licenciar o Windows, a empresa contrata um fornecedor identificável, de quem pode exigir prazo de resposta e correção, em contrapartida, as licenças Open Source fornecem o software no estado em que se encontra, e excluem garantia. Deixando assim o risco com quem adota a ferramenta, salvo em casos de contratação de suporte comercial como o Ubuntu Pro. Isso significa que empresas podem contratar um SLA diretamente com a Microsoft, o que não seria possível (ao menos não no mesmo nível) em soluções Open Source.

\subsection{Comparação em ambientes de servidor}
% TODO: Re-ver esse parágrafo
A vantagem de compatibilidade destacada para o Windows é inerente à sua versão de desktop e se inverte no ambiente de servidores, onde o Linux predomina~\cite{linuxPreferedServer}. Desde 2017, a totalidade dos quinhentos supercomputadores mais rápidos do mundo executa Linux~\cite{top500Computers}, e a própria Microsoft em 2019 informou que maior parte das máquinas virtuais hospedadas no Azure já utilizava o sistema~\cite{azureUsesLinux}. Concorrem para esse quadro o licenciamento por núcleo, que penaliza o crescimento horizontal do parque, o menor consumo de recursos de uma instalação sem interface gráfica, a administração remota por linha de comando e o fato de o software de servidor ser majoritariamente desenvolvido tendo o Linux como alvo primário.

O Windows Server 2025 é licenciado por núcleo físico, com mínimo de oito núcleos por processador e dezesseis por servidor, vendidos em pacotes, sendo US\$ 1.176 pelo pacote de dezesseis núcleos da edição Standard e US\$ 6.771 pela Datacenter~\cite{windowsServerPricing}. Sobre isso ainda incidem as CALs, licenças de acesso exigidas de cada usuário ou dispositivo que autentica no domínio e usa compartilhamento de arquivos~\cite{windowsServerPricing}, e o acesso remoto por RDS cobra uma CAL própria. Este último é de particular destaque, pois significa que o custo cresce também com o número de pessoas na empresa acessando o servidor por RDS, independentemente dos custos do servidor em si.

Outro ponto significativo na operação de um servidor é o reinício para aplicar updates. Atualização do Windows Server requerem reboot para aplicação de mudanças, o que significa uma manutenção programada em serviços sem redundancia ou a redução da capacidade máxima. O Windows Server 2025 trouxe o hotpatching, que aplica a maior parte das correções com a máquina em execução, e mesmo assim mantém cerca de quatro reinícios por ano para as atualizações de baseline~\cite{hotpatchMicrosoft, hotpatchAnuncio}. O recurso exige edição Standard ou Datacenter conectada ao Azure Arc e foi lançado em julho de 2025 cobrando US\$ 1,50 por núcleo de CPU ao mês inicialmente~\cite{hotpatchMicrosoft, hotpatchAnuncio}. A Microsoft ouviu as críticas da comunidade e zerou a cobrança em 15 de maio de 2026~\cite{hotpatchGratuito}, mas manteve a exigência do Azure Arc.

No Ubuntu o equivalente é o Livepatch, que corrige o kernel em execução e vem incluso no Ubuntu Pro, gratuito em até cinco máquinas de uso pessoal e US\$ 500 por ano por servidor no plano pago~\cite{ubuntuProPricing}. Este Livepatch, porém, não cobre todas as atualizações, somente as vulnerabilidades de kernel classificadas como alta ou crítica, e não envolve as atualizações em glibc, OpenSSL, systemd ou microcódigo, que continuam exigindo atualização de pacote e reinício~\cite{livepatchEscopo}. Todavia, isso significa que a máquina reinicia no momento que a equipe escolher, não quando o CVE for publicado.

Todavia, no caso do Windows, é importante destacar sua funcionalidade de controlador de domínio, o Active Directory~\cite{activeDirectory} e o Group Policy~\cite{groupPolicy} não possuem uma ferramenta com maturidade comparavel no ecossistema aberto. Esta função é de extrema importância para escritórios corporativos rodando múltiplos desktops, e é um ponto de destaque do Windows.


\subsection{Pontos Positivos e Negativos:}
\subsubsection{Windows:}
Positivos:
\begin{itemize}
	\item Por ser o SO mais utilizado, é compatível com quase todos os aplicativos desktop utilizados no mundo, desde ferramentas profissionais até jogos;
	\item Suporte profissional direto do fabricante, incluindo suporte 24 horas em planos corporativos;
	\item Facilidade de uso: o design e layout do Windows é aquele ao qual a maioria dos profissionais está acostumada, funcionando como padrão de UI usabilidade;
	\item Active Directory e Group Policy para ambientes corporativos;
	\item Melhor suporte a drivers e hardware, pois fabricantes costumam priorizar compatibilidade com Windows.
\end{itemize}

Negativos:
\begin{itemize}
	\item Custo da licença;
	\item Segurança por obscuridade;
	\item Bloatware;
	\item Obsolescência de hardware imposta por novas versões;
	\item Falta de transparência;
	\item Vulnerabilidades mais conhecidas e exploradas, por ser o SO mais utilizado;
	\item Atualizações forçadas e coleta de telemetria, com pouco controle do usuário sobre esses processos;
	\item Ciclo vicioso de compatibilidade em torno do ecossistema Microsoft.
\end{itemize}

\subsubsection{Ubuntu:}
Positivos:
\begin{itemize}
	\item Não há custo de licença;
	\item Transparência e código auditável pela comunidade;
	\item Comunidade colaborativa, que troca constantemente informações, dados e guias;
	\item Maior segurança em partes, pois a ausência de segurança por obscuridade favorece a descoberta e correção de falhas pela própria comunidade, muito embora nem sempre isso ocorra (vide mais detalhes na Q3);
	\item Flexibilidade para modificar o código-fonte e adequar o sistema à realidade de cada usuário ou empresa;
	\item Redução da dependência de uma única empresa fornecedora;
	\item Suporte LTS longo e menor exigência de hardware;
	\item Disponibilidade de suporte profissional pago para empresas que precisam de SLA formal, como o Ubuntu Pro da Canonical, o que reduz a distância em relação ao suporte oferecido por soluções proprietárias, embora em uma escala menor se comparado à Microsoft.
\end{itemize}

Negativos:
\begin{itemize}
	\item Menor compatibilidade com aplicativos de desktop, por ter parcela do mercado consumidor muito menor que do Windows;
	\item Suporte a hardware e drivers mais limitado, sobretudo para periféricos e componentes recentes;
	\item Falhas de segurança conhecidas publicamente enquanto ainda estão sendo corrigidas\footnote{Embora a descoberta destas vulnerabilidades possa ser oculta durante o processo de reparo, o código vulnerável em si é público. Isso permite a descoberta do erro, assim como é possível chamar a atenção a essas falhas por meio dos commits de correção. Estas publicações permitem a chamada n-day vulnerability, na qual se aproveita o tempo entre o patch e os usuários efetivamente aplicarem as atualizações para realizar o exploit de uma falha, esta janela é geralmente considerada entre 60 e 150 dias~\cite{nDay}} além de outras formas de patch gap~\cite{patchGap};
	\item Fragmentação entre distribuições, o que pode dificultar a escolha e a padronização dentro de uma organização;
	\item Necessidade de treinamento da equipe, caso não haja familiaridade com o sistema.
\end{itemize}

Por tais razões, na utilização de desktops corporativos o Windows continua sendo a escolha principal, pela compatibilidade de aplicativos e drivers e pela usabilidade. Para servidores e infraestrutura o argumento se inverte, e a escolha é o Ubuntu Server por motivos de preço, compatibilidade, e eficiência.